\subsection{Generación automática de rúbricas}
En relación con la construcción automática de estos criterios, Li y sus colaboradores propusieron ARES, un método que transforma documentos en preguntas, respuestas de referencia y rúbricas específicas para cada ejemplo, con el objetivo de generar datos para el entrenamiento por refuerzo en tareas abiertas \cite{li2026ares}. Para ello, el proceso incorpora distintas perspectivas de usuario y características del dominio durante la generación de las preguntas, además de controles para comprobar que puedan responderse sin acceder al documento original, que las respuestas estén respaldadas por su contenido y que las rúbricas cubran los aspectos relevantes de cada consulta. A partir de este procedimiento, los autores construyeron un conjunto de 101.847 ejemplos y 1.108.163 criterios distribuidos en diez dominios, que utilizaron para estudiar el entrenamiento de Qwen3-4B y comparar distintas señales de recompensa, entre ellas las rúbricas específicas, las rúbricas generales y las respuestas de referencia. En cuanto a los resultados, la variante ARES-RL obtuvo el mejor promedio en las siete evaluaciones consideradas, aunque esa ventaja no se mantuvo en todos los escenarios, ya que en la comparación de recompensas sobre HealthBench alcanzó un puntaje de 41,45, frente a 45,97 con rúbricas generales y 50,13 con respuestas de referencia \cite{li2026ares}. El artículo identifica esta evaluación como HealthBench, pero no especifica si corresponde al conjunto original, a HealthBench Hard o a HealthBench Consensus, y señala que las tareas abiertas se evalúan mediante los protocolos de rúbricas o jueces propios de cada benchmark, sin detallar el modelo evaluador de HealthBench ni la fórmula de agregación utilizada para obtener su puntaje global. Por ello, no puede comprobarse a partir de esa descripción si reproduce íntegramente el cálculo oficial. De esta forma, la investigación de ARES muestra cómo generar rúbricas a partir de documentos, pero sus resultados también indican que los criterios específicos de cada consulta no siempre ofrecen la mejor señal de entrenamiento.

